\documentclass[12pt,a4paper]{article}
\usepackage[utf8]{inputenc}
\usepackage[vietnamese]{babel}
\usepackage{amsmath,amssymb,amsthm}
\usepackage{geometry}
\usepackage{enumitem}
\usepackage[most]{tcolorbox}
\usepackage{hyperref}
\usepackage{fancyhdr}
\usepackage{tikz}
\usepackage{eso-pic}
\usepackage{xcolor}
\geometry{a4paper, margin=2cm, bottom=2.5cm}
\hypersetup{colorlinks=true, linkcolor=blue!70!black}
\setlist[itemize]{topsep=4pt, itemsep=2pt}
\setlist[enumerate]{topsep=4pt, itemsep=2pt}
\AddToShipoutPictureFG{%
  \AtPageCenter{%
    \begin{tikzpicture}[remember picture, overlay]
      \node[rotate=45, scale=1, opacity=0.06]{\fontsize{60}{72}\selectfont\bfseries\sffamily Đặng Tấn Quí};
    \end{tikzpicture}%
  }%
}
\pagestyle{fancy}
\fancyhf{}
\fancyhead[L]{\small\textit{THỐNG KÊ ĐA BIẾN}}
\fancyhead[R]{\small\textit{Đặng Tấn Quí}}
\fancyfoot[C]{\thepage}
\fancyfoot[L]{\footnotesize\textcopyright\ Đặng Tấn Quí}
\fancyfoot[R]{\footnotesize SĐT: 0377\,122\,210}
\renewcommand{\headrulewidth}{0.4pt}
\renewcommand{\footrulewidth}{0.4pt}
\fancypagestyle{plain}{\fancyhf{}\fancyfoot[C]{\thepage}\fancyfoot[L]{\footnotesize\textcopyright\ Đặng Tấn Quí}\fancyfoot[R]{\footnotesize SĐT: 0377\,122\,210}\renewcommand{\headrulewidth}{0pt}\renewcommand{\footrulewidth}{0.4pt}}
\newtcolorbox{dinhnghia}[1][]{enhanced, breakable, colback=blue!3!white, colframe=blue!60!black, fonttitle=\bfseries, title={Định nghĩa}, #1}
\newtcolorbox{dinhly}[1][]{enhanced, breakable, colback=green!3!white, colframe=green!50!black, fonttitle=\bfseries, title={Định lý}, #1}
\newtcolorbox{tinhchat}[1][]{enhanced, breakable, colback=yellow!5!white, colframe=orange!50!black, fonttitle=\bfseries, title={Tính chất}, #1}
\newtcolorbox{phuongphap}[1][]{enhanced, breakable, colback=gray!5!white, colframe=gray!50!black, fonttitle=\bfseries, title={Phương pháp}, #1}
\newtcolorbox{luuy}[1][]{enhanced, breakable, colback=red!3!white, colframe=red!50!black, fonttitle=\bfseries, title={Lưu ý}, #1}
\newtcolorbox{congthuc}[1][]{enhanced, breakable, colback=cyan!3!white, colframe=cyan!50!black, fonttitle=\bfseries, title={Công thức}, #1}
\newtcolorbox{vidu}[1][]{enhanced, breakable, colback=violet!3!white, colframe=violet!50!black, fonttitle=\bfseries, title={Ví dụ}, #1}

\begin{document}
\thispagestyle{empty}
\begin{tikzpicture}[remember picture, overlay]
  \fill[blue!80!black] (current page.south west) rectangle (current page.north east);
  \node[anchor=west, white, font=\fontsize{26}{32}\bfseries\selectfont]
    at ([xshift=3cm, yshift=-7.5cm]current page.north west) {THỐNG KÊ ĐA BIẾN};
  \node[anchor=west, white, font=\fontsize{18}{24}\selectfont]
    at ([xshift=3cm, yshift=-9cm]current page.north west) {PCA, phân tích nhân tố, phân biệt, phân cụm};
  \node[anchor=west, white, font=\Large\bfseries] at ([xshift=3cm, yshift=-13cm]current page.north west) {Biên soạn:};
  \node[anchor=west, yellow!80!white, font=\fontsize{22}{28}\bfseries\selectfont] at ([xshift=3cm, yshift=-14.5cm]current page.north west) {ĐẶNG TẤN QUÍ};
  \node[anchor=south east, white, font=\Large\bfseries] at ([xshift=-2cm, yshift=3cm]current page.south east) {\the\year};
\end{tikzpicture}
\newpage
\tableofcontents
\newpage
%% ================================================================
%%  CHƯƠNG 1: PHÂN PHỐI CHUẨN NHIỀU CHIỀU
%% ================================================================
\section{Phân phối chuẩn nhiều chiều}

\begin{dinhnghia}
  $\mathbf{X} = (X_1, \ldots, X_p)^T \sim \mathcal{N}_p(\boldsymbol{\mu}, \boldsymbol{\Sigma})$:
  \[
    f(\mathbf{x}) = \frac{1}{(2\pi)^{p/2}|\boldsymbol{\Sigma}|^{1/2}}\exp\!\left(-\frac{1}{2}(\mathbf{x}-\boldsymbol{\mu})^T\boldsymbol{\Sigma}^{-1}(\mathbf{x}-\boldsymbol{\mu})\right).
  \]
\end{dinhnghia}

\begin{tinhchat}
  \begin{itemize}
    \item Mọi tổ hợp tuyến tính $\mathbf{a}^T\mathbf{X} \sim \mathcal{N}(\mathbf{a}^T\boldsymbol{\mu}, \mathbf{a}^T\boldsymbol{\Sigma}\mathbf{a})$.
    \item Phân phối biên: $X_i \sim \mathcal{N}(\mu_i, \sigma_{ii})$.
    \item Không tương quan $\Leftrightarrow$ độc lập (đặc trưng của chuẩn).
    \item $(\mathbf{X}-\boldsymbol{\mu})^T\boldsymbol{\Sigma}^{-1}(\mathbf{X}-\boldsymbol{\mu}) \sim \chi^2(p)$.
  \end{itemize}
\end{tinhchat}

\begin{dinhly}[title={Phân phối Wishart}]
  Nếu $\mathbf{X}_i \sim \mathcal{N}_p(\boldsymbol{\mu}, \boldsymbol{\Sigma})$ iid thì:
  $\mathbf{W} = \sum_{i=1}^n (\mathbf{X}_i - \bar{\mathbf{X}})(\mathbf{X}_i - \bar{\mathbf{X}})^T \sim W_p(n-1, \boldsymbol{\Sigma})$.
\end{dinhly}

\begin{dinhly}[title={Hotelling $T^2$}]
  $T^2 = n(\bar{\mathbf{X}} - \boldsymbol{\mu}_0)^T \mathbf{S}^{-1}(\bar{\mathbf{X}} - \boldsymbol{\mu}_0)$. Dưới $H_0: \boldsymbol{\mu} = \boldsymbol{\mu}_0$:
  \[
    \frac{n - p}{(n-1)p}\,T^2 \sim F(p, n-p).
  \]
\end{dinhly}

%% ================================================================
%%  CHƯƠNG 2: PHÂN TÍCH THÀNH PHẦN CHÍNH (PCA)
%% ================================================================
\section{Phân tích thành phần chính (PCA)}

\begin{dinhnghia}
  PCA tìm các tổ hợp tuyến tính $Y_k = \mathbf{a}_k^T \mathbf{X}$ sao cho $\text{Var}(Y_k)$ lớn nhất, các $Y_k$ không tương quan.
\end{dinhnghia}

\begin{dinhly}
  Thành phần chính thứ $k$ ứng với \textbf{trị riêng} $\lambda_k$ và \textbf{vectơ riêng} $\mathbf{a}_k$ của $\boldsymbol{\Sigma}$ (hoặc $\mathbf{S}$), $\lambda_1 \ge \lambda_2 \ge \cdots \ge \lambda_p \ge 0$:
  \[
    \text{Var}(Y_k) = \lambda_k, \qquad \text{Tỉ lệ giải thích} = \frac{\lambda_k}{\sum_{j=1}^p \lambda_j}.
  \]
\end{dinhly}

\begin{phuongphap}[title={Thực hành PCA}]
  \begin{enumerate}
    \item Chuẩn hóa dữ liệu (trung bình 0, phương sai 1) nếu các biến khác đơn vị.
    \item Chéo hóa ma trận hiệp phương sai (hoặc tương quan) mẫu.
    \item Chọn số thành phần: scree plot, tỉ lệ phương sai tích lũy $\ge 80\%$.
    \item Giải thích: xem loading (hệ số trong $\mathbf{a}_k$) để hiểu ý nghĩa thành phần.
  \end{enumerate}
\end{phuongphap}

%% ================================================================
%%  CHƯƠNG 3: PHÂN TÍCH NHÂN TỐ (FA)
%% ================================================================
\section{Phân tích nhân tố (Factor Analysis)}

\begin{dinhnghia}
  Mô hình: $\mathbf{X} - \boldsymbol{\mu} = \mathbf{L}\mathbf{F} + \boldsymbol{\varepsilon}$, với $\mathbf{F}$ ($m$ nhân tố chung), $\boldsymbol{\varepsilon}$ (nhiễu riêng), $\mathbf{L}$ (ma trận tải $p \times m$).

  $\boldsymbol{\Sigma} = \mathbf{L}\mathbf{L}^T + \boldsymbol{\Psi}$, $\boldsymbol{\Psi} = \text{diag}(\psi_1, \ldots, \psi_p)$.
\end{dinhnghia}

\begin{phuongphap}[title={Ước lượng và xoay}]
  \begin{itemize}
    \item Ước lượng $\mathbf{L}$: phương pháp thành phần chính, hoặc MLE.
    \item \textbf{Xoay nhân tố}: Varimax (trực giao), Promax (xiên) --- dễ giải thích hơn.
    \item \textbf{Điểm nhân tố}: ước lượng $\hat{\mathbf{F}}$ cho từng quan sát.
  \end{itemize}
\end{phuongphap}

%% ================================================================
%%  CHƯƠNG 4: PHÂN TÍCH PHÂN BIỆT VÀ PHÂN CỤM
%% ================================================================
\section{Phân tích phân biệt và phân cụm}

\subsection{Phân tích phân biệt (Discriminant Analysis)}

\begin{dinhnghia}[title={LDA (Linear Discriminant Analysis)}]
  Phân loại $\mathbf{x}$ vào nhóm $k$ cực đại hóa:
  \[
    \delta_k(\mathbf{x}) = \mathbf{x}^T\boldsymbol{\Sigma}^{-1}\boldsymbol{\mu}_k - \frac{1}{2}\boldsymbol{\mu}_k^T\boldsymbol{\Sigma}^{-1}\boldsymbol{\mu}_k + \ln\pi_k,
  \]
  với $\pi_k = P(\text{nhóm } k)$. Giả sử $\boldsymbol{\Sigma}$ chung.
\end{dinhnghia}

\begin{dinhnghia}[title={QDA (Quadratic DA)}]
  Nới lỏng: mỗi nhóm có $\boldsymbol{\Sigma}_k$ riêng. Biên phân loại là bậc hai.
\end{dinhnghia}

\subsection{Phân cụm (Clustering)}

\begin{phuongphap}[title={K-means}]
  \begin{enumerate}
    \item Chọn $K$ tâm ban đầu.
    \item Gán mỗi điểm vào cụm gần nhất.
    \item Cập nhật tâm = trung bình cụm.
    \item Lặp đến hội tụ. Cực tiểu $\sum_{k=1}^K \sum_{i \in C_k} \|\mathbf{x}_i - \boldsymbol{\mu}_k\|^2$.
  \end{enumerate}
\end{phuongphap}

\begin{phuongphap}[title={Phân cụm phân cấp}]
  \begin{itemize}
    \item \textbf{Agglomerative} (bottom-up): ghép dần.
    \item \textbf{Divisive} (top-down): tách dần.
    \item Liên kết: single, complete, average, Ward.
    \item Kết quả: dendrogram.
  \end{itemize}
\end{phuongphap}

\begin{phuongphap}[title={Chọn số cụm}]
  \begin{itemize}
    \item Phương pháp khuỷu tay (elbow method).
    \item Silhouette score.
    \item Gap statistic.
  \end{itemize}
\end{phuongphap}

\end{document}